\documentclass{book}
\begin{document}
\frontmatter
\title{Notes on Chapter-Level Bibliography Discussion}
\author{Corpus Fixture Working Group}
\date{}
\maketitle
\tableofcontents

\mainmatter
\chapter{Background}
Editorial reports often cite foundational tools at the chapter level before
moving to more recent implementation notes. In this opening chapter we mention
\cite{knuth84}, \cite{lamport94}, and \cite{mittelbach04} because they define the
baseline expectations for layout engines and authoring workflows.

\section{Motivation}
Project teams also reuse algorithmic sources such as \cite{tarjan72} when they
describe dependency ordering for compilation passes.

\chapter{Chapter-Focused Commentary}
The second chapter discusses how a single bibliography can still support
per-chapter narrative structure. Operational guidance from \cite{brooks75} is
often paired with design rules from \cite{hoare69} in retrospective summaries.

\section{Observed Pattern}
The text cites sources from both chapters, which helps exercise the same global
bibliography across multiple chapter headings.

\bibliography{refs}
\end{document}
